\section{Related Work}
\label{sec:related_work}

We review prior work across three key areas: e-commerce chatbots, goal-oriented and persuasive recommendation systems, and the use of knowledge graphs in recommender systems.


% This work builds upon three key research areas: chatbots in e-commerce, goal-oriented and proactive persuasive recommendation systems, and knowledge graphs in recommendation systems.

AI-driven chatbots have been widely adopted in e-commerce to enhance user experience and streamline customer support. \citeauthor{sidlauskiene2023ai}~\cite{sidlauskiene2023ai} showed that using human-like language in AI-driven chatbots can significantly enhance perceived product personalization, with situational loneliness moderating this effect. They also found that the interaction between these factors can increase consumers’ willingness to pay a higher price in conversational commerce. \citeauthor{jiang2022ai}~\cite{jiang2022ai} introduced theoretical models to explain how dialogic chatbot-user interactions affect satisfaction and relationship building. \citeauthor{cheng2020ai}~\cite{cheng2020ai} applied the uses and gratification theory to analyze trade-offs between utility and privacy concerns in chatbot usage. 

However, most existing systems are reactive. They respond to user queries without taking the initiative to steer conversations toward business goals or suggest follow-up actions. This limitation highlights the need for more proactive and goal-oriented conversational agents.

% The integration of AI-powered chatbots into e-commerce platforms has gained significant attention for enhancing customer experience. \citeauthor{sidlauskiene2023ai}~\cite{sidlauskiene2023ai} demonstrated that anthropomorphic design cues significantly influence consumer behavior, with situational loneliness serving as a crucial moderator in user acceptance. \citeauthor{jiang2022ai}~\cite{jiang2022ai} established theoretical frameworks for customer-chatbot relationship dynamics through dialogic interactions that impact satisfaction and engagement. \citeauthor{cheng2020ai}~\cite{cheng2020ai} examined multifaceted factors affecting user experience, applying uses and gratification theory to balance benefits with perceived privacy risks. Despite these advances, most e-commerce chatbots remain reactive systems without actively guiding conversations toward business objectives.

Recent research has increasingly focused on proactive systems that guide users rather than merely respond to their preferences. \citeauthor{bi2024proactive}~\cite{bi2024proactive} proposed the iterative preference guidance (IPG) model, which directs users toward new interests through strategically sequenced recommendations and explicitly defined guiding objectives, helping to mitigate filter bubbles. \citeauthor{jeunen2024multi}~\cite{jeunen2024multi} developed a multivariate policy learning approach for multi-objective recommendation, enabling the joint optimization of competing goals such as user satisfaction and business outcomes. \citeauthor{alslaity2019towards}~\cite{alslaity2019towards} introduced PerPer, a personalized persuasive recommender that incorporates ethical considerations when influencing user decisions. \citeauthor{li2024multi}~\cite{li2024multi} further demonstrated that goal-oriented systems can be deployed effectively even in resource-constrained environments, particularly in educational domains.

% Recent research has shifted toward systems that actively guide users rather than passively responding to preferences. \citeauthor{bi2024proactive}~\cite{bi2024proactive} introduced the iterative preference guidance (IPG) framework, which actively steers users toward new interests through strategically modulated recommendation sequences, addressing filter bubbles by explicitly modeling guiding objectives. \citeauthor{jeunen2024multi}~\cite{jeunen2024multi} advanced multi-objective recommendation through multivariate policy learning, enabling simultaneous optimization of conflicting objectives such as user engagement and business revenue. \citeauthor{alslaity2019towards}~\cite{alslaity2019towards} pioneered persuasive recommendation systems with their personalized persuasive recommender system (PerPer) framework, establishing ethical frameworks for influencing user behavior through personalized approaches. \citeauthor{li2024multi}~\cite{li2024multi} demonstrated practical applications in educational contexts, illustrating the feasibility of implementing goal-oriented systems in resource-constrained environments.

Knowledge graph integration has become a powerful method for modeling structured relationships among entities in recommendation systems. \citeauthor{he2020lightgcn}~\cite{he2020lightgcn} introduced LightGCN, showing that streamlined architectures focusing solely on neighborhood aggregation can outperform more complex graph-based models. \citeauthor{wu2021self}~\cite{wu2021self} applied self-supervised learning techniques to alleviate noise and degree bias in user-item interactions, making the approach particularly effective for platforms with sparse data. \citeauthor{wang2019kgat}~\cite{wang2019kgat} proposed KGAT, which leverages attention mechanisms to model high-order relationships, yielding both accurate and interpretable recommendations. Complementing these efforts, \citeauthor{wang2019neural}~\cite{wang2019neural} developed NGCF, a foundational model that explicitly captures high-order connectivity within user-item interaction graphs using neural networks.


% Knowledge graph integration has emerged as a powerful approach for creating structured relationships between entities. \citeauthor{he2020lightgcn}~\cite{he2020lightgcn} revolutionized graph-based recommendations with LightGCN, demonstrating that simplified architectures focusing on neighborhood aggregation significantly outperform complex models. \citeauthor{wu2021self}~\cite{wu2021self} introduced self-supervised learning to address noisy interactions and degree bias, particularly benefiting enterprises with limited historical data. \citeauthor{wang2019kgat}~\cite{wang2019kgat} introduced KGAT, which models high-order connectivities using attention mechanisms to provide accurate and explainable recommendations. \citeauthor{wang2019neural}~\cite{wang2019neural} established foundational work with neural graph collaborative filtering (NGCF), integrating user-item interaction graphs through explicit high-order connectivity modeling.

Our proposed system integrates key advances from prior research by combining user-friendly chatbot interactions with goal-oriented question suggestion, leveraging knowledge graphs to construct coherent dialogue flows, and employing a multi-agent architecture powered by large language models.



% Our proposed framework synthesizes these research areas by integrating anthropomorphic chatbot design with goal-oriented question suggestion capabilities, leveraging knowledge graphs to create logical conversation pathways while employing multi-agent LLM architectures.